% Generated draft: requires mathematical and source review.
\documentclass[11pt]{article}
\usepackage[margin=25mm]{geometry}
\usepackage{fontspec}
\setmainfont{Latin Modern Roman}
\usepackage{amsmath,amssymb,mathtools,unicode-math}
\setmathfont{Latin Modern Math}
\usepackage{enumitem,needspace,xurl,hyperref}
\setlength{\parindent}{0pt}
\setlength{\parskip}{5pt}
\setlength{\emergencystretch}{3em}
\newcommand{\N}{\mathbb{N}}
\newcommand{\Z}{\mathbb{Z}}
\newcommand{\Q}{\mathbb{Q}}
\newcommand{\R}{\mathbb{R}}
\newcommand{\C}{\mathbb{C}}
\newcommand{\F}{\mathbb{F}}
\newcommand{\A}{\mathbb{A}}
\newcommand{\PP}{\mathbb{P}}
\newcommand{\E}{\mathbb{E}}
\newcommand{\eps}{\varepsilon}
\newcommand{\dd}{\,\mathrm{d}}
\begin{document}
\section*{UnsolvedMath ID: 30000450}
\textbf{Source problem code:} catalogue entry
\subsection{Definitions and mathematical statement}
Fix the root datum and representation-theoretic setting intended in the source. Let \(p\) be its prime characteristic, let \(h\) be the relevant Coxeter number, and assume \(p>h\). Let \(X^+\) denote the dominant integral weights and \(X^+_{\mathrm{reg}}\) the source's regular dominant weights. In the usual affine \(p\)-regular convention, regularity excludes the affine walls
\[
\langle\lambda+\rho,\alpha^\vee\rangle\in p\Z
\]
for roots \(\alpha\), where \(\rho\) is half the sum of positive roots. The source does not reproduce the root datum, algebraic group, base field, or regularity convention, so this conventional interpretation needs confirmation.

Let \(\Delta_{\mathrm{red}}(\lambda)\) and \(\nabla_{\mathrm{red}}(\lambda)\) be the intended reduced standard and costandard modules, let \(l(\lambda)\) be the source's length function, and let \(E^L,E^R\) be its designated subcategories of the ambient derived category. Conjecture 3.2, as identified in the supplied assessment, is
\[
\forall\lambda\in X^+_{\mathrm{reg}},\qquad
\Delta_{\mathrm{red}}(\lambda)[-l(\lambda)]\in E^L
\quad\text{and}\quad
\nabla_{\mathrm{red}}(\lambda)[-l(\lambda)]\in E^R.
\]
The shift is the displayed negative derived-category shift, and the same weight length occurs in both clauses. Both memberships are required for every regular dominant weight, not just weights in the Jantzen region.

The definitions of reduced modules, their integral or quantum-group reduction data, the length convention, the ambient category, and the classes \(E^L,E^R\) are not supplied. They cannot be replaced silently by ordinary Weyl or induced modules, an arbitrary parity convention, or guessed Ext-vanishing conditions. Nor does the workshop title by itself identify the objects as representations of a finite group. These source definitions are necessary for a fully self-contained assertion; the display preserves their precise roles while recording that blocker.

The supplied literature triage checked 2026-08-21 reports conditional and special cases beyond the Jantzen-region discussion, not a proof of the full all-regular-weights assertion. No additional character-formula assumption or weight restriction is added to the conjecture from those partial results.
\subsection{Short English statement}
Under \(p>h\), the two reduced modules, shifted by minus the weight length, are conjectured to lie in the stated left and right categories for every regular dominant weight. The specialized module, length, and category definitions are missing.
\subsection{Sources}
\begin{itemize}
\item Local UnsolvedMath registry, permanent ID 30000450, catalogue entry, \emph{Reduced Standard Modules for Regular Dominant Weights}.
\item \emph{Representations of Finite Groups} (2007), Oberwolfach Reports, pp. 895--978, DOI \url{https://doi.org/10.4171/owr/2006/15}; the supplied assessment identifies Conjecture 3.2.
\item Parshall--Scott, supplied label \emph{Beyond the Jantzen Region}: \url{https://ems.press/content/serial-article-files/46045}.
\item Parshall--Scott, supplied label ``bounding Ext'': \url{https://doi.org/10.1016/j.aim.2010.09.021}.
\item Supplied literature triage checked 2026-08-21; no full specialized definitions or conditional-result hypotheses are reproduced.
\end{itemize}
\end{document}
